\documentclass[11pt,a4paper]{article}

\usepackage[margin=25mm]{geometry}
\usepackage{amsmath}
\usepackage{amssymb}
\usepackage{booktabs}
\usepackage{enumitem}
\usepackage{hyperref}

\title{Self-Written Sequential Linear Programming Method}
\author{DMS4 Optimization Notes}
\date{\today}

\begin{document}

\maketitle

\section{Purpose}

This document describes the optimization method used in the self-written
Sequential Linear Programming (SLP) solver. The purpose of the implementation is
not to reproduce every feature of PySLSQP, but to expose the main numerical
ideas clearly:

\begin{itemize}[nosep]
    \item finite-difference sensitivity estimation,
    \item variable scaling,
    \item linearized constrained subproblems,
    \item slack variables for local feasibility,
    \item a merit function for nonlinear step acceptance,
    \item inexact backtracking line search,
    \item adaptive move limits,
    \item adaptive infeasibility penalty,
    \item convergence checks based on feasibility and step size.
\end{itemize}

The implemented method solves nonlinear constrained optimization problems of the
form
\begin{align}
    \min_{x \in \mathbb{R}^n} \quad & f(x) \\
    \text{subject to} \quad & c_i(x) \geq 0, \qquad i = 1,\dots,m, \\
    & x^L \leq x \leq x^U.
\end{align}

Here, \(f(x)\) is the objective function, typically structural mass, and
\(c_i(x)\) are inequality constraints written in the PySLSQP convention:
positive means feasible and negative means violated.

\section{Notation}

\begin{table}[h]
\centering
\begin{tabular}{ll}
\toprule
Symbol & Meaning \\
\midrule
\(x\) & physical design-variable vector \\
\(x^L, x^U\) & lower and upper bounds on physical variables \\
\(y\) & scaled design-variable vector \\
\(f(x)\) & objective function \\
\(c(x)\) & vector of nonlinear inequality constraints \\
\(g_x = \nabla_x f(x)\) & objective gradient with respect to physical variables \\
\(J_x = \nabla_x c(x)\) & constraint Jacobian with respect to physical variables \\
\(g_y = \nabla_y f(x(y))\) & objective gradient with respect to scaled variables \\
\(J_y = \nabla_y c(x(y))\) & constraint Jacobian with respect to scaled variables \\
\(\Delta y\) & SLP step in scaled variables \\
\(s\) & nonnegative LP slack variables \\
\(M\) & infeasibility/slack penalty parameter \\
\(\delta\) & scaled move-limit vector \\
\(\alpha\) & line-search step length \\
\bottomrule
\end{tabular}
\end{table}

\section{Variable Scaling}

The design variables have different physical units and magnitudes. For example,
a radius may be around \(200\) mm, while a wall thickness may be around \(2\) mm.
Solving the LP directly in the physical variables can make the subproblem poorly
scaled. Therefore, the SLP subproblem is solved in scaled variables.

For variables with finite bounds, the scaling is
\begin{align}
    y_j &= \frac{x_j - x^L_j}{x^U_j - x^L_j}, \\
    x_j &= x^L_j + (x^U_j - x^L_j)y_j.
\end{align}

For variables without a finite range, a fallback scale is used:
\begin{align}
    \text{scale}_j = \max(1, |x_j|).
\end{align}

In compact form,
\begin{align}
    x = x^{\text{off}} + D y,
\end{align}
where \(D\) is a diagonal matrix of variable scales.

Gradients are transformed by the chain rule:
\begin{align}
    g_y &= D g_x, \\
    J_y &= J_x D.
\end{align}

The LP step is computed in \(y\)-space, then converted back to the physical
design space before evaluating the true nonlinear model.

\section{Finite-Difference Sensitivities}

The solver estimates gradients by finite differences. The self-written driver
uses a PySLSQP-like relative finite-difference convention:
\begin{align}
    h_j = r \max(1, |x_j|),
\end{align}
where \(r\) is the relative finite-difference step, normally \(10^{-3}\).

For a forward difference, the objective derivative is approximated as
\begin{align}
    \frac{\partial f}{\partial x_j}
    \approx
    \frac{f(x + h_j e_j) - f(x)}{h_j}.
\end{align}

The constraint derivative is approximated as
\begin{align}
    \frac{\partial c_i}{\partial x_j}
    \approx
    \frac{c_i(x + h_j e_j) - c_i(x)}{h_j}.
\end{align}

Near bounds, the method avoids evaluating outside the allowed box. If a forward
step would violate an upper bound, a backward step is used instead:
\begin{align}
    \frac{\partial f}{\partial x_j}
    \approx
    \frac{f(x) - f(x - h_j e_j)}{h_j}.
\end{align}

If central differences are requested and both forward and backward perturbations
fit within the bounds, the derivative is
\begin{align}
    \frac{\partial f}{\partial x_j}
    \approx
    \frac{f(x + h_j e_j) - f(x - h_j e_j)}{2h_j}.
\end{align}

The bound-aware rule is important for structural design variables such as wall
thicknesses, where evaluating outside the allowed interval may create invalid
geometry or misleading derivatives.

\section{Constraint Violation and Merit Function}

The inequality convention is \(c_i(x) \geq 0\). The violation of one constraint
is therefore
\begin{align}
    v_i(x) = \max(0, -c_i(x)).
\end{align}

The total violation is the \(L_1\) sum
\begin{align}
    V(x) = \sum_{i=1}^{m} v_i(x)
         = \sum_{i=1}^{m} \max(0, -c_i(x)).
\end{align}

The merit function combines objective value and nonlinear infeasibility:
\begin{align}
    \Phi(x; M) = f(x) + M V(x).
\end{align}

Here, \(M > 0\) is the infeasibility penalty. A step is accepted primarily based
on whether it reduces this merit function sufficiently. The merit function is
the bridge between the constrained nonlinear problem and the scalar acceptance
test used by the line search.

\section{SLP Linearization}

At iteration \(k\), the nonlinear objective and constraints are approximated
locally around the current design \(x_k\), or equivalently \(y_k\).

The first-order objective model is
\begin{align}
    f(x(y_k + \Delta y))
    \approx
    f_k + g_{y,k}^T \Delta y.
\end{align}

The first-order constraint model is
\begin{align}
    c(x(y_k + \Delta y))
    \approx
    c_k + J_{y,k} \Delta y.
\end{align}

The linearization is only locally accurate. This is why the method also needs
move limits and backtracking. The LP step is a proposal, not a step that can be
trusted unconditionally.

\section{Move Limits}

The SLP step is restricted by scaled move limits:
\begin{align}
    -\delta_j \leq \Delta y_j \leq \delta_j.
\end{align}

The global variable bounds also restrict the step:
\begin{align}
    y^L_j - y_{k,j} \leq \Delta y_j \leq y^U_j - y_{k,j}.
\end{align}

The final LP bounds for each design variable are
\begin{align}
    \max(y^L_j - y_{k,j}, -\delta_j)
    \leq
    \Delta y_j
    \leq
    \min(y^U_j - y_{k,j}, \delta_j).
\end{align}

The move limit acts as a simple trust region. It prevents the LP from using the
linear model too far away from the current nonlinear evaluation point.

\section{Slack-Relaxed LP Subproblem}

The LP subproblem includes nonnegative slack variables \(s_i\). The slacks relax
the linearized constraints so that the LP remains feasible even if the current
linearization cannot satisfy every constraint exactly.

The subproblem is
\begin{align}
    \min_{\Delta y, s} \quad
        & g_{y,k}^T \Delta y + M \sum_{i=1}^{m} s_i \\
    \text{subject to} \quad
        & c_k + J_{y,k}\Delta y + s \geq 0, \\
        & s_i \geq 0, \qquad i = 1,\dots,m, \\
        & \max(y^L - y_k, -\delta)
          \leq \Delta y \leq
          \min(y^U - y_k, \delta).
\end{align}

The equivalent standard upper-bound inequality form used by linear programming
solvers is
\begin{align}
    -J_{y,k}\Delta y - s \leq c_k.
\end{align}

The slack penalty \(M\sum s_i\) discourages artificial feasibility. Large slacks
indicate that the LP cannot satisfy the linearized constraints within the
current move limits, or that the penalty is too weak compared with the objective
gradient.

\section{Predicted Merit Decrease}

The LP produces a proposed step \(\Delta y\). Before evaluating the nonlinear
model, the solver estimates the merit value predicted by the linear model.

The predicted objective is
\begin{align}
    f_k^{\text{lin}}(\Delta y) = f_k + g_{y,k}^T \Delta y.
\end{align}

The predicted constraints are
\begin{align}
    c_k^{\text{lin}}(\Delta y) = c_k + J_{y,k}\Delta y.
\end{align}

The predicted merit is
\begin{align}
    \Phi_k^{\text{lin}}(\Delta y; M)
    =
    f_k^{\text{lin}}(\Delta y)
    +
    M \sum_{i=1}^{m}
    \max\left(0, -c_{k,i}^{\text{lin}}(\Delta y)\right).
\end{align}

The predicted decrease is
\begin{align}
    \Delta \Phi_{\text{pred}}
    =
    \max\left(
        0,
        \Phi(x_k; M) - \Phi_k^{\text{lin}}(\Delta y; M)
    \right).
\end{align}

This quantity measures what the LP expected to gain.

\section{Inexact Line Search}

The solver uses an inexact backtracking line search along the LP step. The LP
gives a direction \(\Delta y\), and the line search tests
\begin{align}
    y_{\text{trial}} = y_k + \alpha \Delta y,
\end{align}
with
\begin{align}
    \alpha \in \{1, \beta, \beta^2, \beta^3,\dots\},
\end{align}
where \(0 < \beta < 1\). The default value is \(\beta = 0.5\).

The physical trial design is
\begin{align}
    x_{\text{trial}} = x^{\text{off}} + D y_{\text{trial}},
\end{align}
clipped to the physical bounds as a final safeguard.

\subsection{Armijo-like Merit Acceptance}

The main acceptance rule is an Armijo-like sufficient decrease condition:
\begin{align}
    \Phi(x_{\text{trial}}; M)
    \leq
    \Phi(x_k; M)
    -
    \eta \Delta \Phi_{\text{pred}}(\alpha),
\end{align}
where \(\eta\) is a small sufficient-decrease parameter, typically \(10^{-4}\).

This is called inexact because the line search does not try to minimize the
merit function exactly along the direction. It only asks for enough improvement
to justify accepting the step.

\subsection{Feasibility-Progress Acceptance}

When the current design is infeasible, the solver may also accept a step that
significantly reduces total nonlinear violation, even if the merit condition is
not ideal. The condition is
\begin{align}
    V(x_{\text{trial}})
    \leq
    (1-\gamma)V(x_k),
\end{align}
where \(\gamma\) is a small positive reduction fraction.

This protects the method from rejecting useful feasibility-restoration steps
just because the objective temporarily worsens.

\section{Actual-vs-Predicted Quality Ratio}

After evaluating the nonlinear model at the trial point, the actual merit
decrease is
\begin{align}
    \Delta \Phi_{\text{act}}
    =
    \Phi(x_k; M) - \Phi(x_{\text{trial}}; M).
\end{align}

The quality ratio is
\begin{align}
    \rho =
    \frac{\Delta \Phi_{\text{act}}}
         {\Delta \Phi_{\text{pred}}}.
\end{align}

The interpretation is:
\begin{itemize}[nosep]
    \item \(\rho \approx 1\): the linear model predicted the nonlinear response well,
    \item \(\rho > 0\): the step improved the merit function,
    \item \(\rho \leq 0\): the step failed to improve the true nonlinear merit,
    \item small \(\rho\): the linear model was too optimistic.
\end{itemize}

The solver records \(\rho\) in the iteration history. This is one of the most
important diagnostics for understanding why an SLP method rejects steps.

\section{Adaptive Move-Limit Update}

The move limits are updated using the quality ratio:
\begin{align}
    \delta_{k+1}
    =
    \begin{cases}
        \min(a_{\text{exp}}\delta_k, \delta_{\max}),
        & \rho > 0.75, \\
        a_{\text{shr}}\delta_k,
        & \rho < 0.25, \\
        \delta_k,
        & \text{otherwise}.
    \end{cases}
\end{align}

Here,
\begin{align}
    a_{\text{exp}} > 1, \qquad 0 < a_{\text{shr}} < 1.
\end{align}

The idea is simple:
\begin{itemize}[nosep]
    \item if the linear model is reliable, allow larger future steps,
    \item if the linear model is unreliable, reduce the local trust region,
    \item if the model is acceptable but not excellent, keep the same move limit.
\end{itemize}

If the LP fails or no backtracked trial step is acceptable, the move limits are
also reduced.

\section{Adaptive Penalty Update}

The same penalty \(M\) is used in both the LP slack penalty and the nonlinear
merit function. A fixed value of \(M\) can be too weak or too strong depending
on the scaling of the objective and constraints.

The implementation increases \(M\) when constraint slacks or nonlinear
infeasibility persist:
\begin{align}
    M_{k+1} = \min(a_M M_k, M_{\max}),
\end{align}
where \(a_M > 1\).

The purpose is to make future LP subproblems and merit tests place more emphasis
on feasibility when the optimizer keeps trying to improve the objective while
leaving constraints unresolved.

\section{Convergence Criteria}

The solver does not treat a collapsed move limit alone as successful
convergence. Successful convergence requires nonlinear feasibility and a small
design step.

The main success condition is
\begin{align}
    \max_i v_i(x_k) &\leq \varepsilon_{\text{feas}}, \\
    \max_i s_i &\leq \varepsilon_{\text{slack}}, \\
    \|x_k - x_{k-1}\|_2 &\leq \varepsilon_x.
\end{align}

A secondary success condition allows convergence when the design is feasible,
the objective change is small, and the scaled gradient is small:
\begin{align}
    \max_i v_i(x_k) &\leq \varepsilon_{\text{feas}}, \\
    \max_i s_i &\leq \varepsilon_{\text{slack}}, \\
    |f_k - f_{k-1}| &\leq \varepsilon_f, \\
    \|g_{y,k}\|_\infty &\leq \varepsilon_g.
\end{align}

If repeated rejections shrink the move limits below the minimum allowed size,
the solver stops with a failure message. This means the globalization strategy
could not find a reliable step, not that the optimum was necessarily reached.

\section{Complete Algorithm}

At iteration \(k\), the method is:

\begin{enumerate}
    \item Evaluate \(f(x_k)\) and \(c(x_k)\).
    \item Estimate \(g_x\) and \(J_x\) using bound-aware finite differences.
    \item Transform sensitivities to scaled coordinates:
    \[
        g_y = Dg_x, \qquad J_y = J_xD.
    \]
    \item Build and solve the slack-relaxed LP subproblem for \(\Delta y\).
    \item Compute the LP model's predicted merit decrease.
    \item Try backtracking step lengths
    \[
        \alpha = 1, \beta, \beta^2, \dots
    \]
    until the nonlinear trial point satisfies sufficient merit decrease or
    feasibility progress.
    \item If accepted, update \(x\), \(y\), objective, constraints, and history.
    \item Update the move limit using \(\rho\).
    \item Increase the penalty \(M\) if slacks or infeasibility persist.
    \item Check convergence.
    \item If no acceptable step is found, shrink the move limit and try again.
\end{enumerate}

\section{How the Components Connect}

The method components are connected as follows:

\begin{itemize}
    \item Finite differences provide the local sensitivities \(g_x\) and \(J_x\).
    \item Scaling converts those sensitivities to \(g_y\) and \(J_y\), improving
    the numerical balance of the LP.
    \item The LP uses the scaled linear model to propose a constrained search
    direction \(\Delta y\).
    \item Slacks keep the LP solvable and reveal when the linearized constraints
    cannot be satisfied within the current move limit.
    \item The merit function converts objective and constraint violation into a
    scalar value for accepting or rejecting nonlinear trial points.
    \item Backtracking prevents the method from blindly taking a full LP step
    when the nonlinear model disagrees with the linearized model.
    \item The ratio \(\rho\) measures the quality of the linear model and drives
    the move-limit update.
    \item The adaptive penalty increases pressure toward feasibility when
    constraints remain problematic.
    \item The convergence criteria require actual nonlinear feasibility, not
    merely LP feasibility.
\end{itemize}

\section{Important Practical Notes}

\subsection{SLP Versus SLSQP}

SLP solves a sequence of linear programming subproblems. SLSQP solves a sequence
of quadratic programming subproblems, where the quadratic model contains
curvature information. Because the self-written method uses only a linear model,
it relies heavily on scaling, move limits, and line search to avoid poor steps.

\subsection{Why Rejections Happen}

A rejection means the LP's linearized prediction did not produce an acceptable
improvement when checked on the true nonlinear functions. Common causes are:

\begin{itemize}[nosep]
    \item move limits are too large,
    \item finite-difference gradients are noisy,
    \item active constraints are highly nonlinear,
    \item the penalty \(M\) is too low,
    \item the LP is dominated by objective reduction instead of feasibility,
    \item the current design is close to a bound and derivatives are one-sided.
\end{itemize}

\subsection{Useful Diagnostics}

The most useful recorded quantities are:

\begin{itemize}[nosep]
    \item \(\alpha\): accepted line-search step length,
    \item \(\rho\): actual-vs-predicted merit quality,
    \item \(V(x)\): total nonlinear constraint violation,
    \item \(\max_i v_i(x)\): worst nonlinear constraint violation,
    \item \(\sum_i s_i\): total LP slack,
    \item \(\max_i s_i\): worst LP slack,
    \item \(\max_j \delta_j\): current scaled move-limit size,
    \item \(M\): current merit/slack penalty.
\end{itemize}

If \(\rho\) is often low or negative, the linear model is too optimistic and the
move limits or finite-difference steps should be reviewed. If slacks remain
large, the constraints are difficult to satisfy locally or the penalty is too
weak.

\section{Default Parameter Interpretation}

\begin{table}[h]
\centering
\begin{tabular}{lll}
\toprule
Parameter & Typical value & Interpretation \\
\midrule
\(r\) & \(10^{-3}\) & relative finite-difference step \\
\(M_0\) & \(10^3\) & initial infeasibility/slack penalty \\
\(\delta_{\max}\) & \(0.10\) & maximum scaled move limit \\
\(\delta_{\min}\) & \(10^{-5}\) & minimum scaled move limit before failure stop \\
\(\beta\) & \(0.5\) & backtracking shrink factor \\
\(\eta\) & \(10^{-4}\) & sufficient-decrease fraction \\
\(\varepsilon_{\text{feas}}\) & \(10^{-5}\) & nonlinear feasibility tolerance \\
\(\varepsilon_{\text{slack}}\) & \(10^{-8}\) & LP slack tolerance \\
\(\varepsilon_x\) & \(10^{-4}\) & physical step-size tolerance \\
\bottomrule
\end{tabular}
\end{table}

\section{Summary}

The self-written optimizer is a scaled, slack-relaxed SLP method with
merit-based globalization. The LP subproblem proposes a locally constrained
direction. The nonlinear merit function and backtracking line search decide how
much of that direction can be trusted. The actual-vs-predicted ratio updates the
move limits, while the adaptive penalty increases focus on feasibility when
needed. These pieces are necessary because a linearized LP is only a local
approximation of the true nonlinear structural optimization problem.

\end{document}
